\section{現代テスト理論その2}
\subsection{授業内容}
\subsubsection{科目の中でのこのコマの位置づけ}
項目反応理論の導入によって，被験者母数と項目母数が完全に分離され，項目の特徴を細かく記述することができるようになった。また，項目の特徴がわかっていればテストの実践方法も変わってくる。ひとつはCATに代表されるように，ダイナミックに出題を変化させることができるようになること，そうしたうえでもテストの平均点が事前にコントロールしうることなどが示される。項目から得られる情報という観点から項目情報曲線が，項目情報曲線の累積からテスト情報曲線が導出される。
つづいてこのテスト理論の発展形として，多段階モデルに拡張可能なことをみる。特に段階反応モデルは，リッカートのシグマ法のように段階反応をモデルかできるという意味で，現代的リッカート法であるともいえる。段階反応モデルを用いることで，適切な反応段階のチェックをすることができるなど，応用的側面が高いことを確認する。
またテスト理論は因子分析の特殊系であるという扱いだったが，多段階，多因子へと展開することで再び因子分析モデルに統合されていくことを確認する。

% \include{lessonII05}現代テスト理論；信頼性の新しい考え方・段階反応理論


\subsubsection{コマ主題細目}
\begin{description}
    \item[現代テスト理論の特徴]
    \item[情報曲線]
    \item[段階反応モデル] 
    \item[因子分析の歴史と展開]  
\end{description}
\subsubsection{キーワード}
\begin{itemize}
\item 通過率
\item ロジスティックモデル
\item 被験者母数の推定について
\end{itemize}
\subsection{授業情報}
\paragraph{コマの展開方法}
講義
\subsubsection{予習・復習課題}
\paragraph{予習}
\paragraph{復習}
